% Part: lambda-calculus
% Chapter: lambda-definability
% Section: lists

\documentclass[../../../include/open-logic-section]{subfiles}

\begin{document}

\olfileid{lam}{ldf}{lst}
\olsection{Dhaptar}

Dhaptar $[M_1, M_2, \ldots, M_n]$ bisa ditetepaké dadi:
\[
  [M_1, M_2, \ldots, M_n] \ident \lambd[sz][s M_1 (s M_2 (\ldots (s M_n z)))]
\]
Kanthi cara informal, suku kang merepresentasikaké dhaptar
iku fungsi ``panglumpuk'' dhaptar kasebut: fungsi kang nampa
rong suku. Suku kapisan minangka fungsi kang nampa nilai kang
wis diklumpukaké lan sawijining unsur anyar, banjur mbalèkaké
nilai panglumpuk anyar. Suku kapindho minangka nilai wiwitan.
Fungsi iki nglumpukaké unsur siji-siji lan pungkasané
mbalèkaké asilé.

\begin{digress}
  Yèn wis kulina karo pemrograman fungsional, definisi iki
  mirip karo \emph{lipetan tengen}: sawijining dhaptar
  ditetepaké minangka fungsi lipetan tengen saka unsur-unsuré.
\end{digress}

Saka enkoding iki, kita bisa kanthi gampang netepaké
sawetara fungsi kang migunani:
\begin{align*}
  \fn{Sum} & \ident
  \lambd[l][l \, (\lambd[xa][\fn{Add} \, x \, a])\,  \num{0}]\\
  \fn{Len} & \ident
  \lambd[l][l \, (\lambd[xa][\fn{Add} \, a \, \num{1}]) \num{0}]
\end{align*}

$\fn{Sum}$ ngitung gunggung saka sawijining dhaptar numeral
Church. Fungsi iki nglumpukaké nilai ing dhaptar, kanthi nilai
wiwitan $\num{0}$, lan kanggo saben unsur mbalèkaké
$\fn{Add} \, x\, a$ minangka nilai anyar (ing kéné $x$
iku unsur saiki lan $a$ nilai kang wis diklumpukaké).
Asilé yaiku gunggung kabèh unsur.

\begin{prob}
  Apa kang ditindakake $\fn{Len}$? Terangna.
\end{prob}

\begin{prob}
  Tetepna fungsi kang mbalikké urutan sawijining dhaptar.
\end{prob}
\end{document}
